% ============================================================================
%  Kripta Studios — Feature Space Specification
%  GBT+RL 0DTE Options Trading System
% ============================================================================
\documentclass[11pt,a4paper,twoside]{article}

\usepackage[utf8]{inputenc}
\usepackage[T1]{fontenc}
\usepackage{lmodern}
\usepackage{amsmath, amsfonts, amssymb, amsthm}
\usepackage{mathtools}
\usepackage{bm}
\usepackage{geometry}
\geometry{margin=1in}
\usepackage{titlesec}
\usepackage{setspace}
\onehalfspacing
\usepackage{graphicx}
\usepackage{booktabs}
\usepackage{multirow}
\usepackage{longtable}
\usepackage{float}
\usepackage{hyperref}
\hypersetup{colorlinks=true, linkcolor=blue, citecolor=red, urlcolor=blue}
\usepackage{listings}
\lstset{basicstyle=\ttfamily\small, breaklines=true, frame=single, columns=fullflexible}
\usepackage{xcolor}

\newtheorem{definition}{Definition}
\newtheorem{proposition}{Proposition}
\newtheorem{remark}{Remark}

\DeclareMathOperator*{\argmax}{arg\,max}
\newcommand{\E}{\mathbb{E}}
\newcommand{\R}{\mathbb{R}}
\newcommand{\bx}{\bm{x}}
\newcommand{\bs}{\bm{s}}
\newcommand{\ba}{\bm{a}}

% ============================================================================
\title{\textbf{Feature Space Specification} \\[0.5em]
\large{163-Dimensional State Representation for the GBT Directional Classifier \\
and PPO Reinforcement Learning Agent}}
\author{Kripta Studios — Quantitative Research Division}
\date{March 2026}

\begin{document}
\maketitle
\tableofcontents
\newpage

% ============================================================================
\section{Introduction}
\label{sec:intro}
% ============================================================================

The directional GBT classifier and the PPO reinforcement learning agent both operate on a shared 163-dimensional feature vector $\bx_t \in \R^{163}$, computed at each 5-minute bar $t$ during regular trading hours (09:30--16:00 ET). This vector encodes the complete market microstructure state as seen from the perspective of a 0DTE options market participant: spot price dynamics, the full Greek surface, dealer inventory exposures, volatility term structure, cross-asset context, and temporal position within the trading session.

The feature set is defined in \texttt{hybrid\_model.py} under the constant \texttt{FEATURE\_COLUMNS}, which is imported identically by the GBT training pipeline, the RL episode generator, and the live inference bot. This guarantees perfect train/test/production parity.

\begin{definition}[Feature Vector]
The feature vector at time $t$ is the ordered tuple:
\begin{equation}
    \bx_t = (x_t^{(1)}, x_t^{(2)}, \ldots, x_t^{(163)}) \in \R^{163}
    \label{eq:feature_vector}
\end{equation}
where each component $x_t^{(i)}$ is a real-valued statistic computed from publicly observable market data up to and including time $t$. No lookahead is permitted: $x_t^{(i)} \perp \{S_\tau : \tau > t\}$.
\end{definition}

% ============================================================================
\section{Feature Taxonomy}
\label{sec:taxonomy}
% ============================================================================

The 163 features are organized into ten functional categories, summarized in Table~\ref{tab:feature_summary}.

\begin{longtable}{llr}
\toprule
\textbf{Category} & \textbf{Primary Information Content} & \textbf{Dimensions} \\
\midrule
\endhead
\midrule
\endfoot
\bottomrule
\caption{Summary of feature categories and their dimensionality.}
\label{tab:feature_summary}
\endlastfoot
Spot \& Log Returns         & Price level, 1/5/15/30-min returns, vol-adj returns & 8 \\
Technical Indicators        & RSI, MACD, Bollinger Bands, momentum & 10 \\
Greek Exposures             & $\Delta$, $\Gamma$, $\theta$, $\mathcal{V}$, $\rho$ by expiry bucket & 24 \\
GEX / DEX / VEX             & Aggregate gamma, delta, vanna exposure measures & 18 \\
Put/Call Metrics            & Volume PCR, OI PCR, delta-filtered PCR & 12 \\
IV Surface                  & ATM IV, skew, term structure, call/put smile slopes & 16 \\
Order Flow Derivatives      & Net gamma, net vanna, net charm, DGEX, cross-products & 14 \\
Level Proximity             & Distance to max-gamma strike, IB levels ($D{-}1$ to $D{-}5$) & 21 \\
Temporal Encoding           & Time-of-day sinusoid, day-of-week sinusoid, minutes-to-close & 5 \\
Cross-Ticker Context        & QQQ/VIX/TLT returns, spread, correlation & 35 \\
\midrule
\textbf{Total} & & \textbf{163} \\
\end{longtable}

% ============================================================================
\section{Category Definitions}
\label{sec:categories}
% ============================================================================

\subsection{Spot Price and Log Returns ($D = 8$)}

The raw spot price $S_t$ is standardized relative to its recent distribution:
\begin{equation}
    x^{\text{spot}}_t = \frac{S_t - \bar{S}_{20}}{s_{20}}, \quad \bar{S}_{20} = \frac{1}{20}\sum_{i=1}^{20} S_{t-i},\quad s_{20} = \text{std}\{S_{t-i}\}_{i=1}^{20}
    \label{eq:spot_norm}
\end{equation}

Multi-horizon log returns are included at four lookback windows $\Delta \in \{1, 5, 15, 30\}$ minutes:
\begin{equation}
    r_{t,\Delta} = \log\left(\frac{S_t}{S_{t-\Delta}}\right)
    \label{eq:log_return}
\end{equation}

Volatility-adjusted counterparts normalize each return by the implied 0DTE ATM volatility $\sigma^{\text{ATM}}_t$ scaled to the corresponding horizon:
\begin{equation}
    \tilde{r}_{t,\Delta} = \frac{r_{t,\Delta}}{\sigma^{\text{ATM}}_t \sqrt{\Delta / (252 \times 390)}}
    \label{eq:vol_adj_return}
\end{equation}
This transformation maps raw price dynamics into a stationary, approximately standard-normal space, conditioning the GBT on moves expressed in implied-volatility units rather than raw basis points.

\subsection{Technical Indicators ($D = 10$)}

Classical price-based indicators are included as non-linear compressors of recent price history:
\begin{itemize}
    \item \textbf{RSI(14):} 14-period Relative Strength Index, bounded in $[0, 100]$.
    \item \textbf{MACD and signal:} 12-period minus 26-period exponential moving averages, plus the 9-period signal line and histogram.
    \item \textbf{Bollinger Bands:} upper band $\mu_{20} + 2\sigma_{20}$, lower band $\mu_{20} - 2\sigma_{20}$, and bandwidth $4\sigma_{20} / \mu_{20}$.
\end{itemize}

\subsection{Greek Exposures ($D = 24$)}

Aggregate net Greeks across the listed options chain are computed for both 0DTE and the nearest weekly expiration:
\begin{equation}
    \mathcal{G}^{(k)}_t = \sum_{\substack{i \in \text{chain}\\T_i = T_k}} \text{OI}_i \cdot \text{Greek}_i, \quad k \in \{\text{0DTE}, \text{weekly}\}
    \label{eq:net_greek}
\end{equation}
where $\text{OI}_i$ is the open interest of contract $i$. Six Greeks ($\Delta, \Gamma, \theta, \mathcal{V}, \rho$, charm) are computed for each of the two expiry buckets, yielding 12 features per side (24 total).

\subsection{GEX, DEX, and VEX ($D = 18$)}
\label{subsec:gex_dex}

Dealer exposure measures estimate the aggregate hedging pressure that market makers exert on the spot price:

\begin{definition}[Gamma Exposure (GEX)]
The market-maker net gamma exposure at time $t$ is:
\begin{equation}
    \text{GEX}_t = \sum_{i \in \text{chain}} \Gamma_i \cdot \text{OI}_i \cdot \text{Multiplier} \cdot S_t^2 / 100
    \label{eq:gex}
\end{equation}
A positive $\text{GEX}_t$ indicates dealers are net long gamma: they sell rallies and buy dips, imposing mean-reverting pressure on the spot price. Negative GEX indicates dealers are net short gamma, amplifying directional moves.
\end{definition}

\begin{definition}[Delta Exposure (DEX) and Vanna Exposure (VEX)]
\begin{align}
    \text{DEX}_t &= \sum_{i} \Delta_i \cdot \text{OI}_i \cdot \text{Multiplier} \cdot S_t \\
    \text{VEX}_t &= \sum_{i} \mathcal{V}^{(\text{vanna})}_i \cdot \text{OI}_i \cdot \text{Multiplier}
\end{align}
DEX measures the directional delta hedge flow. VEX captures the sensitivity of dealer delta to changes in implied volatility, becoming particularly relevant during VIX spike regimes.
\end{definition}

These three exposure metrics are computed separately for 0DTE, weekly, and the combined chain, yielding 18 features.

\subsection{Put/Call Metrics ($D = 12$)}

Put/call ratios provide an aggregate measure of options market sentiment:
\begin{align}
    \text{PCR}_{\text{vol}} &= \frac{V_{\text{put}}}{V_{\text{call}}} \\
    \text{PCR}_{\text{OI}}  &= \frac{\text{OI}_{\text{put}}}{\text{OI}_{\text{call}}} \\
    \text{PCR}_{\Delta}     &= \frac{\sum_{i:\Delta_i < 0} |\Delta_i| \cdot \text{OI}_i}{\sum_{j:\Delta_j > 0} \Delta_j \cdot \text{OI}_j}
\end{align}
where $V$ denotes traded volume. The delta-filtered PCR weights contracts by their hedge ratio, providing a more precise measure of directional pressure than raw volume ratios. These are computed for three expiry tiers (0DTE, weekly, all).

\subsection{Implied Volatility Surface ($D = 16$)}

The IV surface features encode the term structure and cross-sectional shape of implied volatility:
\begin{itemize}
    \item \textbf{ATM IV (0DTE, weekly, monthly):} implied volatility at the at-the-money strike.
    \item \textbf{IV Skew:} $\text{IV}_{25\Delta,\text{put}} - \text{IV}_{25\Delta,\text{call}}$, a risk-reversal proxy for demand asymmetry.
    \item \textbf{Term structure:} weekly/0DTE IV ratio and monthly/weekly IV ratio.
    \item \textbf{Smile slopes:} OTM call slope and OTM put slope, capturing tail-risk premium asymmetry.
\end{itemize}

\subsection{Order Flow Derivatives ($D = 14$)}

Higher-order Greek cross-products capture the sensitivity of dealer hedge flows to joint price-volatility moves:
\begin{align}
    \text{net\_vanna}_t  &= \sum_i \frac{\partial^2 C_i}{\partial S \, \partial \sigma} \cdot \text{OI}_i \\
    \text{net\_charm}_t  &= \sum_i \frac{\partial^2 C_i}{\partial S \, \partial t} \cdot \text{OI}_i \\
    \text{DGEX}_t        &= \frac{d(\text{GEX})}{dS}\bigg|_t
\end{align}
These terms become significant near zero-day expiration when the term-structure derivatives of gamma approach infinity.

\subsection{Level Proximity and Initial Balance ($D = 21$)}
\label{sec:ib_features}

Market microstructure research has long identified that traders react asymmetrically to the initial price range established during the first hour of trading. The Initial Balance (IB) for day $d$ is:
\begin{equation}
    \text{IB}_{\text{high}}^{(d)} = \sup_{\tau \in [09:30, 10:30]_d} S_\tau, \qquad
    \text{IB}_{\text{low}}^{(d)}  = \inf_{\tau \in [09:30, 10:30]_d} S_\tau
    \label{eq:ib_definition}
\end{equation}

Distance features quantify the spot price's position relative to IB levels for the current and prior five trading days $d \in \{0, -1, -2, -3, -4, -5\}$:
\begin{equation}
    \delta^{(d)}_{\text{high}} = \text{clip}\!\left(\frac{S_t - \text{IB}_{\text{high}}^{(d)}}{\nu \cdot S_t},\; -10,\; 10\right), \quad \nu = 0.005
    \label{eq:ib_distance}
\end{equation}
An inside-IB boolean $\mathbf{1}[{\text{IB}_{\text{low}}^{(0)} \leq S_t \leq \text{IB}_{\text{high}}^{(0)}}]$ is also included.

Additionally, confluence features encode whether the current spot price aligns with the maximum-gamma strike $K^*_\Gamma$ via a Gaussian RBF kernel:
\begin{equation}
    \Phi^{\Gamma}_t = \exp\!\left(-\frac{(S_t - K^*_\Gamma)^2}{2\,(0.005 \cdot S_t)^2}\right), \quad K^*_\Gamma = \argmax_{K} \sum_{i: K_i = K} \Gamma_i \cdot \text{OI}_i
    \label{eq:gamma_confluence}
\end{equation}
$\Phi^{\Gamma}_t \to 1$ when spot approaches the dominant gamma strike; $\Phi^{\Gamma}_t \to 0$ in open price space. Analogous confluence features are computed for Vanna and DGEX strikes.

\subsection{Temporal Encoding ($D = 5$)}

Intraday and weekly periodicity is encoded via Fourier projections to prevent artificial discontinuities at the day and week boundaries that would otherwise confuse tree-based splits:
\begin{align}
    x^{\sin}_{\text{tod}} &= \sin\!\left(2\pi \cdot \frac{m_t}{390}\right), \quad
    x^{\cos}_{\text{tod}} = \cos\!\left(2\pi \cdot \frac{m_t}{390}\right) \\
    x^{\sin}_{\text{dow}} &= \sin\!\left(2\pi \cdot \frac{d_t}{5}\right), \quad
    x^{\cos}_{\text{dow}} = \cos\!\left(2\pi \cdot \frac{d_t}{5}\right)
\end{align}
where $m_t \in [0, 390]$ is the number of minutes elapsed since the 09:30 open, and $d_t \in \{0,1,2,3,4\}$ is the day-of-week index. A raw \texttt{minutes\_to\_close} $= 390 - m_t$ feature is also included to allow the GBT to model the accelerating theta-decay regime in the final hour.

\subsection{Cross-Ticker Context ($D = 35$)}

The system tracks four secondary instruments: QQQ, VIX, TLT, and the SPX/QQQ spread. For each, multi-horizon returns and selected Greeks are included:
\begin{itemize}
    \item \textbf{VIX:} spot level, 1-min and 5-min returns. Acts as an implied-volatility regime indicator.
    \item \textbf{TLT (20Y Treasury ETF):} 1-min, 5-min, and 15-min returns. Captures risk-off flows.
    \item \textbf{QQQ:} full analog of the SPX feature vector (spot, returns, Greeks). Encodes tech-sector divergence.
    \item \textbf{SPX/QQQ spread:} normalized difference in 5-minute returns. Captures sector rotation.
\end{itemize}

% ============================================================================
\section{Numerical Preprocessing Pipeline}
\label{sec:preprocessing}
% ============================================================================

\subsection{Log Transform for Extreme-Scale Features}

Features with raw magnitudes exceeding $10^4$ (e.g., raw GEX in dollar-gamma units) undergo the sign-preserving log transform before being written to the training Parquet:
\begin{equation}
    \tilde{x} = \text{sign}(x) \cdot \log(1 + |x|)
    \label{eq:log_transform}
\end{equation}
This is applied to: \texttt{net\_gamma}, \texttt{net\_vanna}, \texttt{net\_charm}, \texttt{net\_dgex}, \texttt{gex\_0dte}, \texttt{gex\_weekly}, \texttt{delta\_gamma\_ratio}, and all raw exposure scalars.

\subsection{Feature Normalizer}

The \texttt{FeatureNormalizer} class fits per-feature mean $\mu_j$ and standard deviation $\sigma_j$ on the training window:
\begin{equation}
    x_{\text{norm}}^{(j)} = \frac{x^{(j)} - \mu_j}{\sigma_j + \varepsilon}, \quad \varepsilon = 10^{-8}
    \label{eq:normalizer}
\end{equation}
Statistics are fit once per walk-forward window and serialized to \texttt{models/hybrid\_normalizer\_wf.npz}. At inference time, the normalizer is loaded and applied identically to every input batch, including the live trading bot.

\subsection{NaN Handling}

The GBT training pipeline applies \texttt{np.nan\_to\_num(X, nan=0.0, posinf=5.0, neginf=-5.0)} before model fitting. LightGBM natively handles missing values (NaN) via its optimal split-finding mechanism, which learns a dedicated branch direction for NaN samples. The explicit imputation to zero is applied only as a safety measure for features that should never be NaN in well-formed data.

% ============================================================================
\section{Feature Importance and Selection}
\label{sec:importance}
% ============================================================================

LightGBM provides three native feature importance metrics: split count, gain, and permutation importance. Across walk-forward windows, the consistently highest-importance features are:

\begin{table}[H]
\centering
\begin{tabular}{clcc}
\toprule
\textbf{Rank} & \textbf{Feature} & \textbf{Category} & \textbf{Gain Contribution} \\
\midrule
1 & \texttt{vix\_spot}                & Cross-Ticker     & High \\
2 & \texttt{ret\_5m\_vol\_adj}         & Spot \& Returns  & High \\
3 & \texttt{gex\_0dte}                & GEX/DEX/VEX      & Medium-High \\
4 & \texttt{pcr\_vol\_0dte}            & Put/Call         & Medium-High \\
5 & \texttt{iv\_skew\_0dte}            & IV Surface       & Medium-High \\
6 & \texttt{dist\_ib\_high}            & Level Proximity  & Medium \\
7 & \texttt{net\_vanna}               & Order Flow       & Medium \\
8 & \texttt{minutes\_to\_close}        & Temporal         & Medium \\
9 & \texttt{confluence\_max\_gamma}    & Level Proximity  & Medium \\
10 & \texttt{tlt\_ret\_5m}             & Cross-Ticker     & Medium \\
\midrule
\multicolumn{4}{l}{\textit{Note: Rankings are window-averaged; individual windows may vary.}} \\
\bottomrule
\end{tabular}
\caption{Top-10 features by average LightGBM gain across walk-forward windows.}
\label{tab:importance}
\end{table}

\begin{remark}
The dominance of \texttt{vix\_spot} and \texttt{ret\_5m\_vol\_adj} confirms that volatility regime and recent momentum are the primary drivers of 0DTE directional predictability. The GEX and IB features provide secondary conditioning, particularly for mean-reversion signals near dominant gamma strikes.
\end{remark}

% ============================================================================
\section{RL State Space Extension}
\label{sec:rl_state}
% ============================================================================

The PPO agent observes an extended 182-dimensional state vector $\bs_t$ that appends 19 dynamic and position-management features to the base 163-dimensional market feature vector:
\begin{equation}
    \bs_t = \left[\bx_t \;\Big|\; \bx^{\text{dyn}}_t \;\Big|\; \bm{q}_t \;\Big|\; \bm{c}_t \right] \in \R^{182}
    \label{eq:rl_state}
\end{equation}
where:
\begin{itemize}
    \item $\bx^{\text{dyn}}_t \in \R^8$ — Dynamic market features: per-minute spot history, delta history, and volume.
    \item $\bm{q}_t \in \R^7$ — Position state: unrealized PnL\%, hold time, $|\Delta|$, recovery probability, implied volatility, maximum adverse excursion (MAE), and trailing drawdown.
    \item $\bm{c}_t \in \R^4$ — Model context: GBT confidence score $c_t$, time-to-target, dummy edge indicator, and $\log\sigma$ (observation noise).
\end{itemize}

The state is augmented during training via \texttt{augment\_state()}, which adds small Gaussian noise $\mathcal{N}(0, 0.01)$ to all continuous components to improve policy generalization.

\begin{thebibliography}{6}

\bibitem{grinsztajn2022}
L.~Grinsztajn, E.~Oyallon, and G.~Varoquaux,
``Why do tree-based models still outperform deep learning on typical tabular data?''
\textit{NeurIPS}, vol.~35, 2022.

\bibitem{ke2017}
G.~Ke et al.,
``LightGBM: A Highly Efficient Gradient Boosting Decision Tree,''
\textit{NeurIPS}, vol.~30, 2017.

\bibitem{cboe_gex}
S.~Avellaneda and A.~Savine,
``Hedging with options: Gamma exposure and market impact,''
\textit{Journal of Financial Economics}, 2021.

\bibitem{initial_balance}
D.~Steidlmayer,
``Markets and Market Logic,'' Chicago Board of Trade, 1986.

\bibitem{schulman2017}
J.~Schulman et al.,
``Proximal Policy Optimization Algorithms,''
\textit{arXiv:1707.06347}, 2017.

\end{thebibliography}

\end{document}